\section*{Введение}
\addcontentsline{toc}{section}{Введение}

\textbf{Актуальность темы.} В современных автоматизированных системах управления
технологическим процессом (АСУТП) программируемые логические контроллеры (ПЛК)
используются как базовые исполнительные узлы, от которых зависит корректность
работы оборудования, устойчивость технологического режима и безопасность
эксплуатации. Для таких систем критичны не только надежность алгоритмов
управления, но и организованное сопровождение инженерных данных: конфигураций,
параметров, журналов событий, резервных копий и сервисных комментариев.
Рекомендации NIST по защите industrial control systems подчеркивают, что для
промышленных систем существенны требования к доступности, прослеживаемости
изменений и контролируемому доступу к данным~\cite{nist_ics}. Стандарт ISA-95,
в свою очередь, рассматривает интеграцию производственного контура и
информационных систем верхнего уровня как самостоятельную практическую
задачу~\cite{isa95}.

На практике данные о ПЛК нередко хранятся разрозненно: часть сведений находится в
локальных каталогах инженеров, часть --- в сетевых папках, часть --- в архивах
проектов и выгрузках журналов. Такой подход затрудняет поиск актуальной версии
конфигурации, делает неудобным анализ неисправностей и осложняет разграничение
доступа между участниками сопровождения. Дополнительную специфику предметной
области создает то, что инженеры АСУТП часто работают в локальных сетях объектов,
не имеющих постоянного доступа к сети Интернет. Следовательно, практическая
система сопровождения должна учитывать как централизованную веб-работу с общей
базой данных, так и возможность безопасного получения данных из локального
контура без прямого удаленного управления оборудованием.

\textbf{Степень разработанности проблемы.} На рынке присутствуют решения,
закрывающие отдельные части рассматриваемой задачи. CODESYS Automation Server
предоставляет централизованные средства учета устройств, работы с версиями и
управления доступом~\cite{codesys_as}. Siemens TIA Portal Project Server решает
задачи совместной работы с инженерными проектами~\cite{siemens_multiuser}.
GitLab поддерживает хранение версий, историю изменений и организацию командной
работы~\cite{gitlab_platform,gitlab_issues}. Grafana Loki ориентирована на сбор,
хранение и поиск журналов событий~\cite{loki_overview}. При этом для контроллеров
LogicBox существует открытый репозиторий host-side software, содержащий утилиты
для обнаружения устройств и получения инженерных данных в локальной сети,
например \texttt{lbdiscover}, \texttt{lbcmd}, \texttt{lbconf} и
\texttt{lbgetlog}~\cite{logicbox_repo}. Однако указанные средства не образуют
единой веб-системы, ориентированной на учет объектов эксплуатации, накопление
истории конфигураций, разграничение доступа по ролям и областям ответственности,
а также синхронизацию результатов локальной инженерной работы.

\textbf{Объект исследования} --- процессы учета, хранения, поиска и анализа
конфигурационных и диагностических данных контроллеров LogicBox в составе АСУТП.

\textbf{Предмет исследования} --- методы и программные средства построения
веб-приложения для централизованного учета ПЛК LogicBox, хранения связанных
файлов и синхронизации инженерных данных, полученных из локальной сети объекта.

\textbf{Цель работы} --- разработать концепцию веб-приложения, предназначенного
для централизованного учета контроллеров LogicBox, хранения конфигурационных и
диагностических данных, разграничения прав доступа пользователей и приема данных,
собранных локальным вспомогательным клиентом без реализации опасных функций
удаленного управления оборудованием.
